\documentclass{setp-new}
\usepackage{fontspec}
\setmainfont{Times New Roman}
\setsansfont{Arial}
\IfFontExistsTF{Songti SC}{\newcommand{\setpSongFont}{Songti SC}}{\newcommand{\setpSongFont}{SimSun}}
\newif\ifsetpLocalNotoHei
\IfFileExists{/Users/zhouleixishu/Library/Fonts/NotoSansCJKsc-Regular.otf}{%
  \setpLocalNotoHeitrue}{%
  \setpLocalNotoHeifalse}
\ifsetpLocalNotoHei\else
\IfFontExistsTF{SimHei}{\newcommand{\setpHeiFont}{SimHei}}{%
  \IfFontExistsTF{Noto Sans CJK SC}{\newcommand{\setpHeiFont}{Noto Sans CJK SC}}{%
    \PackageError{ReSETP}{No approved Simplified Chinese Heiti font found}{Install SimHei or Noto Sans CJK SC before compiling}%
    \newcommand{\setpHeiFont}{Noto Sans CJK SC}}}
\fi
\IfFontExistsTF{Songti SC}{\newcommand{\setpFangSongFont}{Songti SC}}{\newcommand{\setpFangSongFont}{FangSong}}
\ifsetpLocalNotoHei
\setCJKmainfont[ItalicFont=\setpSongFont,SlantedFont=\setpSongFont]{\setpSongFont}
\setCJKsansfont[Path=/Users/zhouleixishu/Library/Fonts/]{NotoSansCJKsc-Regular.otf}
\setCJKmonofont{\setpSongFont}
\setCJKfamilyfont{zhsong}{\setpSongFont}
\setCJKfamilyfont{zhhei}[Path=/Users/zhouleixishu/Library/Fonts/]{NotoSansCJKsc-Regular.otf}
\setCJKfamilyfont{zhfs}{\setpFangSongFont}
\else
\setCJKmainfont[BoldFont=\setpHeiFont,ItalicFont=\setpSongFont,SlantedFont=\setpSongFont]{\setpSongFont}
\setCJKsansfont{\setpHeiFont}
\setCJKmonofont{\setpSongFont}
\setCJKfamilyfont{zhsong}[BoldFont=\setpHeiFont]{\setpSongFont}
\setCJKfamilyfont{zhhei}{\setpHeiFont}
\setCJKfamilyfont{zhfs}{\setpFangSongFont}
\fi
\providecommand{\songti}{}
\providecommand{\heiti}{}
\providecommand{\fangsong}{}
\renewcommand{\songti}{\CJKfamily{zhsong}}
\renewcommand{\heiti}{\CJKfamily{zhhei}}
\renewcommand{\fangsong}{\CJKfamily{zhfs}}
\XeTeXlinebreaklocale "zh"
\XeTeXlinebreakskip = 0pt plus 1pt
\usepackage{mathtools}
\usepackage{booktabs}
\usepackage{array}
\usepackage{tabularx}
\usepackage{float}
\usepackage{ragged2e}
\newcolumntype{Y}{>{\RaggedRight\arraybackslash}X}
\newcommand{\setpsymtabsetup}{\fontsize{9pt}{10.5pt}\selectfont\renewcommand{\arraystretch}{0.95}\setlength{\tabcolsep}{1.8pt}}
\setlength{\parindent}{2em}
\setlength{\emergencystretch}{3em}
\setlength{\headheight}{32pt}
\setlength{\abovedisplayskip}{3pt plus .6pt minus .6pt}
\setlength{\belowdisplayskip}{3pt plus .6pt minus .6pt}
\setlength{\abovedisplayshortskip}{2pt plus .4pt minus .4pt}
\setlength{\belowdisplayshortskip}{2pt plus .4pt minus .4pt}
\setlength{\jot}{1.5pt}
\raggedbottom
\allowdisplaybreaks
\renewcommand{\tablename}{表}
\jolname{系统工程理论与实践}{Systems Engineering --- Theory \& Practice}
\Volume{2026}{5}{0}{0}
\PageNum{1}
\PaperDOI{}
\Category{U492.3}{A}
\TitleMark{周雷习书等: 时变碳强度下多车场动态协同路径建模与算法}
\Title{时变碳强度下多车场动态协同路径建模与算法}
\Author{周雷习书, 干宏程, 邱莹莹}{上海理工大学 管理学院, 上海 200093}
\Abstract{本文建立时变电网碳强度下多车场动态协同混合车队路径优化模型。}
\Keywords{多车场协同配送; 混合车队; 动态需求; 非线性充电; 时变碳强度}
\ETitle{Multi-depot collaborative routing under time-varying carbon intensity}
\EAuthor{ZHOU Leixishu, GAN Hongcheng, QIU Yingying}{Business School, University of Shanghai for Science and Technology, Shanghai 200093, China}
\EAbstract{A modular formulation is developed for dynamic collaborative multi-depot mixed-fleet routing under time-varying grid carbon intensity.}
\EKeywords{collaborative routing; mixed fleet; dynamic demand; nonlinear charging; carbon intensity}
\begin{document}
\setcounter{section}{1}
\section{模型建立}

\subsection{问题描述}

考虑由多个配送企业组成的协同配送系统。每家企业拥有一个车场及有限的油电混合车队。
在每次决策时刻，已到达客户的需求量和硬时间窗均为已知。协作条件下，车辆可直接服务其他企业负责的客户，
但不在车场之间转运货物。车辆由所属车场出发并返回，可在一个运营日内连续执行多个配送趟；
电动车可在车场或公共充电站部分充电。充电费用和间接碳排放分别由实际充电时段的电价
和电网碳强度确定。模型以运营成本与碳交易成本之和最小为目标，确定车辆使用、
配送路径、趟次衔接和充电安排。

模型假设如下：1）客户需求不可拆分，每个客户仅由一辆车服务一次；
2）车辆载重、客户时间窗和电池容量均为硬约束；
3）同一车辆相邻配送趟在时间上不得重叠；
4）车辆从所属车场完成补货，运营日结束后返回所属车场；
5）电动车允许部分充电，充电功率随电池电量变化；
6）分时电价和时变电网碳强度为外生参数；
7）协作路径确定后再进行成本分摊，分摊结果不改变路径模型。
模型符号说明如表\ref{tab:unified-symbols}所示。

\begin{table}[H]
\centering
\caption{符号说明}
\label{tab:unified-symbols}
\setpsymtabsetup
\begin{tabularx}{\linewidth}{@{}c>{\RaggedRight\arraybackslash}Xc>{\RaggedRight\arraybackslash}X@{}}
\toprule
符号 & 说明 & 符号 & 说明\\
\midrule
$D$ & 车场集合 & $N$ & 客户集合\\
$S$ & 充电站集合 & $V$ & 车场、客户及充电站节点集合\\
$K$ & 车辆集合 & $K^{g}$ & 燃油车集合\\
$K^{e}$ & 电动车集合 & $\mathcal T$ & 电价和碳强度核算时段集合\\
$\Omega_k$ & 车辆$k$的可行配送方案集合 & $\mathcal R_{kp}$ & 方案$p$的配送趟集合\\
$\mathcal C_{kp}$ & 方案$p$的充电过程集合 & $\mathcal C_{kp}^{rr'}$ & 配送趟$r$与$r'$间的充电过程集合\\
$V_r$ & 配送趟$r$的节点集合 & $V_r^{c}$ & 配送趟$r$服务的客户集合\\
$r\prec_p r'$ & 方案$p$中$r'$紧接在$r$之后 & $r_1$ & 方案中的首个配送趟\\
$H$ & 不含车场副本的任意节点子集 & $h$ & 充电过程索引\\
$d_r^{+}$ & 配送趟$r$的起点车场副本 & $d_r^{-}$ & 配送趟$r$的终点车场副本\\
$x_{ijr}$ & 配送趟$r$使用弧$(i,j)$时取1 & $a_{ir}$ & 配送趟$r$访问节点$i$时取1\\
$z_{kp}$ & 车辆$k$采用方案$p$时取1 & $\alpha_{ikp}$ & 方案$p$服务客户$i$时为1，否则为0\\
$\beta_{skpt}$ & 方案$p$在充电地点$s$、时段$t$充电时为1，否则为0 & $C_s$ & 充电地点$s$的充电桩数\\
$q_i$ & 客户$i$的需求量 & $[e_i,l_i]$ & 节点$i$的时间窗\\
$s_i$ & 客户$i$的服务时间 & $d_{ij}$ & 节点$i$至$j$的距离\\
$t_{ij}^{k}$ & 车辆$k$在弧$(i,j)$上的行驶时间 & $Q_k$ & 车辆$k$的最大载重\\
$B_k$ & 电动车$k$的电池容量 & $B_k^{\min}$ & 电动车$k$的最低允许电量\\
$\ell_{ir}$ & 车辆离开节点$i$时的载重 & $\tau_{ir}$ & 车辆在节点$i$的服务开始时刻\\
$w_{ir}$ & 车辆在节点$i$的停留时间 & $\varepsilon_{ir}$ & 车辆到达节点$i$时的电量\\
$Y_{ir}$ & 车辆在节点$i$的补电量 & $v_{ij}^{k}$ & 车辆$k$在弧$(i,j)$上的行驶速度\\
$b_{ij}^{k}$ & 电动车$k$在弧$(i,j)$上的电耗 & $f_{ij}^{k}$ & 燃油车$k$在弧$(i,j)$上的油耗\\
$c_k^{\mathrm f}$ & 车辆$k$的固定成本 & $c_k^{\mathrm m}$ & 车辆$k$的单位里程成本\\
$p^{\mathrm f}$ & 单位燃油价格 & $p^{\mathrm c}$ & 单位碳价格\\
$\overline E$ & 系统碳排放配额 & $D_{kp}$ & 方案$p$的总里程\\
$G_{kp}$ & 方案$p$的总油耗 & $C_{kp}^{\mathrm e}$ & 方案$p$的充电费用\\
$E_{kp}^{g}$ & 方案$p$的燃油车直接排放 & $E_{kp}^{e}$ & 方案$p$的电动车间接排放\\
$\gamma_t$ & 时段$t$的电网碳强度 & $\lambda^{g}$ & 燃油排放因子\\
$\sigma(h)$ & 充电过程$h$所在站点 & $B_h^{a}$ & 充电过程$h$开始时的电量\\
$B_h^{d}$ & 充电过程$h$结束时的电量 & $\Delta_h$ & 充电过程$h$的持续时间\\
$B_{ht}$ & 充电过程$h$在时段$t$端点的电量 & $e_{ht}$ & 充电过程$h$在时段$t$的补电量\\
$p_{st}^{\mathrm e}$ & 站点$s$在时段$t$的单位电价 & $T^{H}$ & 运营时域上限\\
$M$ & 足够大的正数 & $N^{(m)}$ & 第$m$次重规划的待服务客户集合\\
$\Omega_k^{(m)}$ & 第$m$次重规划的可行配送方案集合 & $u(m)$ & 第$m$个批事件处理时刻\\
$u_q^{(m)}$ & 第$m$次定量触发时刻 & $T^{\mathrm{int}}$ & 动态需求处理间隔\\
$q(t_1,t_2)$ & 时段$[t_1,t_2]$内新增需求量 & $\bar q$ & 动态需求累计上限\\
\bottomrule
\end{tabularx}
\end{table}

\subsection{目标函数}

\subsubsection{车辆启动成本}

车辆启动成本主要包括车辆折旧、维修、保险及驾驶员工资等，与投入使用的车辆数成正比。
\begin{equation}
F_1=\sum_{k\in K}\sum_{p\in\Omega_k}c_k^{\mathrm f}z_{kp}.
\label{eq:unified-f1}
\end{equation}

\subsubsection{行驶成本}

行驶成本与车辆行驶里程成正比，主要反映车辆维护、轮胎磨损和通行等非能源费用。
\begin{equation}
F_2=\sum_{k\in K}\sum_{p\in\Omega_k}c_k^{\mathrm m}D_{kp}z_{kp}.
\label{eq:unified-f2}
\end{equation}

\subsubsection{充电成本}

充电成本由电动车的实际补电量及其所处电价时段确定，不同充电时刻对应不同的充电费用。
\begin{equation}
F_3=\sum_{k\in K^{e}}\sum_{p\in\Omega_k}C_{kp}^{\mathrm e}z_{kp}.
\label{eq:unified-f3}
\end{equation}

\subsubsection{油耗成本}

油耗成本与燃油车行驶产生的油耗成正比，车辆总油耗由各配送趟的弧段油耗累计得到。
\begin{equation}
F_4=p^{\mathrm f}\sum_{k\in K^{g}}\sum_{p\in\Omega_k}G_{kp}z_{kp}.
\label{eq:unified-f4}
\end{equation}

\subsubsection{碳排放成本}

碳排放成本由燃油车的直接排放和电动车充电产生的间接排放构成，
并按系统总排放量与碳配额的差额计算。
\begin{equation}
F_5=p^{\mathrm c}\left[
\sum_{k\in K^{g}}\sum_{p\in\Omega_k}E_{kp}^{g}z_{kp}
+\sum_{k\in K^{e}}\sum_{p\in\Omega_k}E_{kp}^{e}z_{kp}
-\overline E
\right].
\label{eq:unified-f5}
\end{equation}

\subsection{能耗、充电与碳排放模型}

\subsubsection{车辆能耗}

采用实际能耗模型计算车辆行驶产生的电耗和油耗\cite{ref:23}。
车辆$k$在弧$(i,j)$上的机械功率、电耗、单位时间油耗率和弧段油耗分别为
\begin{align}
P_{ij}^{k}(\ell_{ir},v_{ij}^{k})
&=\frac{1}{2}c_d\rho^{\mathrm{air}}A_k^{f}(v_{ij}^{k})^3
+(m_k+\ell_{ir})gc_rv_{ij}^{k},
\label{eq:unified-power}\\
b_{ij}^{k}(\ell_{ir},v_{ij}^{k})
&=\phi^{d}\varphi^{d}P_{ij}^{k}(\ell_{ir},v_{ij}^{k})
\frac{d_{ij}}{v_{ij}^{k}},
\label{eq:unified-electricity}\\
g_{ij}^{k}(\ell_{ir},v_{ij}^{k})
&=\xi^{\mathrm f}\frac{\epsilon R^{\mathrm e}
+P_{ij}^{k}(\ell_{ir},v_{ij}^{k})/\eta}{\kappa\psi},
\label{eq:unified-fuel-rate}\\
f_{ij}^{k}(\ell_{ir},v_{ij}^{k})
&=g_{ij}^{k}(\ell_{ir},v_{ij}^{k})\frac{d_{ij}}{v_{ij}^{k}}.
\label{eq:unified-fuel}
\end{align}
式中，$c_d,\rho^{\mathrm{air}},A_k^{f},m_k,g,c_r$分别为空气阻力系数、空气密度、
迎风面积、车辆整备质量、重力加速度和滚动阻力系数；
$\phi^{d},\varphi^{d}$为电驱动参数，
$\xi^{\mathrm f},\epsilon,R^{\mathrm e},\eta,\kappa,\psi$为燃油率模型参数。
$\phi^{d}\varphi^{d}$包含功率至电量的单位换算和驱动效率折算，
$b_{ij}^{k}$和$f_{ij}^{k}$的单位分别为kWh和L。
$D_{kp}$和$G_{kp}$分别由方案$p$所含配送趟的弧段距离和式\eqref{eq:unified-fuel}逐项累计。

\subsubsection{非线性充电与分时电价}

设$b_s^{c}$为站点$s$的充电函数分段数，
$B_{sn}^{c},T_{sn}^{c},\rho_{sn}$分别为第$n$个断点的电量、
累计充电时间和分段斜率。非线性充电函数采用分段线性形式\cite{ref:94}：
\begin{equation}
\Phi_s(\tau)=
\begin{cases}
\rho_{s0}\tau, & \tau\le T_{s0}^{c},\\
B_{s0}^{c}+\rho_{s1}(\tau-T_{s0}^{c}),
& T_{s0}^{c}<\tau\le T_{s1}^{c},\\
\cdots & \cdots\\
B_{s,b_s^{c}-1}^{c}+\rho_{s,b_s^{c}}(\tau-T_{s,b_s^{c}-1}^{c}),
& T_{s,b_s^{c}-1}^{c}<\tau\le T_{s,b_s^{c}}^{c}.
\end{cases}
\label{eq:unified-charge-function}
\end{equation}
设$B_h^{a},B_h^{d},\Delta_h$分别为充电过程$h$开始时的电量、结束时的电量和持续时间，
$\sigma(h)$为充电地点。充电时长为\cite{ref:montoya}
\begin{equation}
\Delta_h=
\Phi_{\sigma(h)}^{-1}(B_h^{d})-
\Phi_{\sigma(h)}^{-1}(B_h^{a}).
\label{eq:unified-charge-duration}
\end{equation}
充电前后的电量满足
\begin{equation}
0\le B_h^{a}\le B_h^{d}\le B_k,
\qquad h\in\mathcal C_{kp},\quad k\in K^{e}.
\label{eq:unified-charge-energy-bound}
\end{equation}

\Needspace{11\baselineskip}
设$U_t$为时段$t$的端点，$p_{st}^{\mathrm e}$为站点$s$在时段$t$的单位电价。
电价与电网碳强度采用统一的时段集合~$\mathcal T$，且$U_0<U_1<\cdots<U_{|\mathcal T|}$\cite{ref:94}，则
\begin{equation}
\Psi_s(\tau)=
\begin{cases}
p_{s1}^{\mathrm e}, & U_0\le\tau\le U_1,\\
\cdots & \cdots\\
p_{st}^{\mathrm e}, & U_{t-1}\le\tau\le U_t,\\
\cdots & \cdots\\
p_{s,|\mathcal T|}^{\mathrm e}, & U_{|\mathcal T|-1}\le\tau\le U_{|\mathcal T|}.
\end{cases}
\label{eq:unified-price-function}
\end{equation}
相邻时段的公共端点归入后一时段。
对于给定充电过程，以$B_{ht}$表示时段$t$端点的电量：充电开始前取$B_h^{a}$，
充电结束后取$B_h^{d}$，充电区间内由式\eqref{eq:unified-charge-function}确定。
时段补电量$e_{ht}$满足\cite{ref:94}
\begin{equation}
e_{ht}=B_{ht}-B_{h,t-1},\qquad
\sum_{t\in\mathcal T}e_{ht}=B_h^{d}-B_h^{a},\qquad
h\in\mathcal C_{kp},\quad k\in K^{e}.
\label{eq:unified-charge-balance}
\end{equation}
方案$p$的充电费用为
\begin{equation}
C_{kp}^{\mathrm e}
=\sum_{h\in\mathcal C_{kp}}\sum_{t\in\mathcal T}
p_{\sigma(h),t}^{\mathrm e}e_{ht}.
\label{eq:unified-charge-cost}
\end{equation}

\subsubsection{时变碳排放}

燃油车排放按油耗量计算，电动车充电间接排放按不同时段的补电量和电网碳强度计算\cite{ref:deng}。
$\lambda^{g}$和$\gamma_t$的单位分别为kgCO$_2$e/L和kgCO$_2$e/kWh：
\begin{align}
E_{kp}^{g}&=\lambda^{g}G_{kp},
\label{eq:unified-direct-emission}\\
E_{kp}^{e}&=\sum_{h\in\mathcal C_{kp}}\sum_{t\in\mathcal T}\gamma_t e_{ht}.
\label{eq:unified-indirect-emission}
\end{align}

\subsection{多车场多趟协同配送模型}

多趟配送按同一车辆在一个运营日内依次执行多个配送趟的方式建立\cite{ref:44}。
对任一方案$p\in\Omega_k$及配送趟$r\in\mathcal R_{kp}$，
$d_r^{+}$和$d_r^{-}$为车辆所属车场的起点和终点副本；
充电站在重复访问时使用相应节点副本。模型的总目标为
\begin{equation}
\min F=F_1+F_2+F_3+F_4+F_5.
\label{eq:unified-objective}
\end{equation}

车辆$k$的可行配送方案集合$\Omega_k$由满足下列约束的配送趟和充电安排构成。
方案属性$D_{kp}$、$G_{kp}$、$C_{kp}^{\mathrm e}$、$E_{kp}^{g}$和$E_{kp}^{e}$以及方案系数
$\alpha_{ikp}$、$\beta_{skpt}$均由相应的配送趟与充电安排确定。
对任一$p\in\Omega_k$及$r\in\mathcal R_{kp}$，为简化表示，以下变量省略下标$k,p$。
车场和充电站节点的需求量取0；$r_1$表示方案$p$的首趟，$r\prec_p r'$表示$r'$紧接在$r$之后。
客户点、充电节点和车场节点的停留时间分别为服务时间、充电时间和趟间作业时间。
每个被访问的公共充电站节点对应一个充电过程$h$，其中
$B_h^{a}=\varepsilon_{ir}$、$B_h^{d}=\varepsilon_{ir}+Y_{ir}$且$w_{ir}=\Delta_h$；
$\mathcal C_{kp}^{rr'}$表示相邻配送趟之间在车场完成的充电过程集合，趟间作业时间包括补货时间及相应充电时间。

\begin{equation}
\begin{aligned}
\sum_{j\in V_r\setminus\{d_r^{+}\}}x_{d_r^{+}jr}&=1, &
\sum_{i\in V_r\setminus\{d_r^{-}\}}x_{id_r^{-}r}&=1 .
\end{aligned}
\label{eq:unified-trip-start-end}
\end{equation}
\begin{equation}
\begin{aligned}
\sum_{i\in V_r}x_{id_r^{+}r}&=0, &
\sum_{j\in V_r}x_{d_r^{-}jr}&=0 .
\end{aligned}
\label{eq:unified-trip-direction}
\end{equation}
\begin{equation}
\sum_{j\in V_r\setminus\{i\}}x_{ijr}
=\sum_{j\in V_r\setminus\{i\}}x_{jir}=a_{ir},
\qquad i\in V_r\setminus\{d_r^{+},d_r^{-}\}.
\label{eq:unified-flow}
\end{equation}
\begin{equation}
\sum_{i\in H}\sum_{\substack{j\in H\\j\ne i}}x_{ijr}
\le |H|-1,
\qquad H\subseteq V_r\setminus\{d_r^{+},d_r^{-}\},\quad |H|\ge2.
\label{eq:unified-subtour}
\end{equation}
\begin{equation}
\begin{aligned}
\ell_{d_r^{+}r}&=\sum_{i\in V_r^{c}}q_i a_{ir}, &
0\le\ell_{ir}&\le Q_k,\qquad i\in V_r .
\end{aligned}
\label{eq:unified-load-bound}
\end{equation}
\begin{equation}
-M(1-x_{ijr})
\le\ell_{jr}-\ell_{ir}+q_j
\le M(1-x_{ijr}),
\qquad i,j\in V_r,\quad i\ne j .
\label{eq:unified-load-propagation}
\end{equation}
\begin{equation}
e_i a_{ir}
\le\tau_{ir}\le l_i+M(1-a_{ir}),
\qquad i\in V_r\setminus\{d_r^{+},d_r^{-}\}.
\label{eq:unified-time-window}
\end{equation}
\begin{equation}
\tau_{jr}
\ge\tau_{ir}+w_{ir}+t_{ij}^{k}-M(1-x_{ijr}),
\qquad i,j\in V_r,\quad i\ne j .
\label{eq:unified-time-propagation}
\end{equation}
\begin{equation}
\tau_{d_{r'}^{+}r'}
\ge\tau_{d_r^{-}r}+w_{d_r^{-}r},
\qquad r\prec_p r' .
\label{eq:unified-trip-link-time}
\end{equation}
\begin{equation}
\tau_{d_r^{-}r}\le T^{H},
\qquad r\in\mathcal R_{kp}.
\label{eq:unified-horizon}
\end{equation}
\begin{equation}
\begin{aligned}
B_k^{\min}\le\varepsilon_{ir}&\le B_k,
&&i\in V_r,\quad k\in K^{e},\\
\varepsilon_{d_{r_1}^{+}r_1}&=B_k,
&&k\in K^{e}.
\end{aligned}
\label{eq:unified-battery-bound}
\end{equation}
\begin{equation}
-M(1-x_{ijr})
\le\varepsilon_{jr}-\varepsilon_{ir}-Y_{ir}+b_{ij}^{k}
\le M(1-x_{ijr}),
\qquad i,j\in V_r,\quad i\ne j,\quad k\in K^{e}.
\label{eq:unified-battery-propagation}
\end{equation}
\begin{equation}
Y_{ir}=0,
\qquad i\in V_r^{c}\cup\{d_r^{+},d_r^{-}\},\quad k\in K^{e}.
\label{eq:unified-no-customer-charge}
\end{equation}
\begin{equation}
0\le Y_{ir}\le B_k-\varepsilon_{ir},
\qquad i\in S\cap V_r,\quad k\in K^{e}.
\label{eq:unified-charge-bound}
\end{equation}
\begin{equation}
\varepsilon_{d_{r'}^{+}r'}
=\varepsilon_{d_r^{-}r}
+\sum_{h\in\mathcal C_{kp}^{rr'}}(B_h^{d}-B_h^{a}),
\qquad r\prec_p r',\quad k\in K^{e}.
\label{eq:unified-trip-link-energy}
\end{equation}
车辆方案选择采用集合划分形式\cite{ref:sahin}。充电设施容量按照同一时段内的充电车辆数确定\cite{ref:12}。
\begin{equation}
\sum_{k\in K}\sum_{p\in\Omega_k}\alpha_{ikp}z_{kp}=1,
\qquad i\in N .
\label{eq:unified-coverage}
\end{equation}
\begin{equation}
\sum_{p\in\Omega_k}z_{kp}\le1,
\qquad k\in K .
\label{eq:unified-vehicle}
\end{equation}
\begin{equation}
\sum_{k\in K}\sum_{p\in\Omega_k}\beta_{skpt}z_{kp}\le C_s,
\qquad s\in S\cup D,\quad t\in\mathcal T .
\label{eq:unified-station}
\end{equation}
\begin{equation}
\begin{aligned}
x_{ijr},a_{ir},z_{kp}&\in\{0,1\},\\
\ell_{ir},\tau_{ir},w_{ir},\varepsilon_{ir},Y_{ir},
B_h^{a},B_h^{d},B_{ht},e_{ht},\Delta_h&\ge0 .
\end{aligned}
\label{eq:unified-domain}
\end{equation}

式\eqref{eq:unified-trip-start-end}和式\eqref{eq:unified-trip-direction}限定每个配送趟的起讫车场及行驶方向；
式\eqref{eq:unified-flow}为被访问节点的流量平衡约束；式\eqref{eq:unified-subtour}用于消除与车场分离的子回路。
式\eqref{eq:unified-load-bound}确定车辆离场时的载重并限制车辆容量，式\eqref{eq:unified-load-propagation}描述服务客户前后的载重变化。
式\eqref{eq:unified-time-window}为节点时间窗约束，式\eqref{eq:unified-time-propagation}为趟内时间连续性约束，
式\eqref{eq:unified-trip-link-time}和式\eqref{eq:unified-horizon}分别限定相邻配送趟的时间关系及车辆返回车场的最迟时刻。

式\eqref{eq:unified-battery-bound}限定电池电量及首趟初始电量；
式\eqref{eq:unified-battery-propagation}描述车辆沿弧行驶和在节点补电后的电量变化；
式\eqref{eq:unified-no-customer-charge}和式\eqref{eq:unified-charge-bound}分别限定允许充电的节点及补电量；
式\eqref{eq:unified-trip-link-energy}保证同一车辆相邻配送趟之间的电量连续。
式\eqref{eq:unified-coverage}为客户唯一服务约束，式\eqref{eq:unified-vehicle}为车辆方案选择约束，
式\eqref{eq:unified-station}为充电设施容量约束，式\eqref{eq:unified-domain}规定决策变量的取值范围。

\subsection{考虑动态需求的协同配送模型}

\subsubsection{目标函数}

第$m$次重规划时，$N^{(m)}$为尚未完成及新到达的客户集合，
$\Omega_k^{(m)}$为保留车辆已执行路径后形成的可行配送方案集合。
上标$(m)$表示第$m$次重规划对应的量。
各成本系数由已执行路径与待规划路径共同确定，动态阶段的目标函数为\cite{ref:71}
\begin{equation}
F_1^{(m)}
=\sum_{k\in K}\sum_{p\in\Omega_k^{(m)}}c_k^{\mathrm f}z_{kp}^{(m)}.
\label{eq:unified-dynamic-f1}
\end{equation}
\begin{equation}
F_2^{(m)}
=\sum_{k\in K}\sum_{p\in\Omega_k^{(m)}}c_k^{\mathrm m}D_{kp}^{(m)}z_{kp}^{(m)}.
\label{eq:unified-dynamic-f2}
\end{equation}
\begin{equation}
F_3^{(m)}
=\sum_{k\in K^{e}}\sum_{p\in\Omega_k^{(m)}}C_{kp}^{\mathrm e,(m)}z_{kp}^{(m)}.
\label{eq:unified-dynamic-f3}
\end{equation}
\begin{equation}
F_4^{(m)}
=p^{\mathrm f}\sum_{k\in K^{g}}\sum_{p\in\Omega_k^{(m)}}G_{kp}^{(m)}z_{kp}^{(m)}.
\label{eq:unified-dynamic-f4}
\end{equation}
\begin{equation}
\begin{aligned}
F_5^{(m)}
&=p^{\mathrm c}\left[
\sum_{k\in K^{g}}\sum_{p\in\Omega_k^{(m)}}E_{kp}^{g,(m)}z_{kp}^{(m)}
+\sum_{k\in K^{e}}\sum_{p\in\Omega_k^{(m)}}E_{kp}^{e,(m)}z_{kp}^{(m)}
-\overline E
\right].
\end{aligned}
\label{eq:unified-dynamic-f5}
\end{equation}
\begin{equation}
\min F^{(m)}
=F_1^{(m)}+F_2^{(m)}+F_3^{(m)}+F_4^{(m)}+F_5^{(m)}.
\label{eq:unified-dynamic-objective}
\end{equation}

\subsubsection{动态信息处理策略}

动态需求采用定时、定量联合批处理策略\cite{ref:71}。
设$[t_0,t_n]$为动态需求接收区间，令$u(0)=t_0$，
$q(t_1,t_2)$为$[t_1,t_2]$内新增需求量之和，则
\begin{align}
u(m)
&=\min\{u_q^{(m)},u(m-1)+T^{\mathrm{int}},t_n\},
\label{eq:unified-trigger-time}\\
u_q^{(m)}
&=\min\{u>u(m-1):q(u(m-1),u)\ge\bar q\}.
\label{eq:unified-trigger-quantity}
\end{align}
到达批事件处理时刻后，保留已执行路径，更新客户集合、车辆可用时刻、载重和电量状态，
并以$N^{(m)}$和$\Omega_k^{(m)}$替换相应集合重新求解。

\begin{thebibliography}{99}
\bibitem{ref:71} 邱莹莹. 低碳背景下考虑动态需求的混合车队物流配送路径问题研究[D]. 上海: 上海理工大学, 2024.\par
\bibitem{ref:23} 陈婉茹, 徐光明, 张得志, 等. 碳交易机制下多中心混合车队配送路径和速度优化研究[J]. 系统工程理论与实践, 2023, 43(11): 3320--3335.\par
\bibitem{ref:94} Wang W, Adulyasak Y, Cordeau J F. Electric vehicle routing with heterogeneous charging stations[J]. Computers \& Operations Research, 2026, 189: 107374.\par
\bibitem{ref:montoya} Montoya A, Guéret C, Mendoza J E, Villegas J G. The electric vehicle routing problem with nonlinear charging function[J]. Transportation Research Part B: Methodological, 2017, 103: 87--110.\par
\bibitem{ref:93} Shi W, Wang N, Zhou L, He Z. The bi-objective mixed-fleet vehicle routing problem under decentralized collaboration and time-of-use prices[J]. Expert Systems with Applications, 2025, 273: 126875.\par
\bibitem{ref:deng} Deng J, Hu H, Gong S, Dai L. Impacts of charging pricing schemes on cost-optimal logistics electric vehicle fleet operation[J]. Transportation Research Part D: Transport and Environment, 2022, 109: 103333.\par
\bibitem{ref:44} Zhen L, Ma C, Wang K, Xiao L, Zhang W. Multi-depot multi-trip vehicle routing problem with time windows and release dates[J]. Transportation Research Part E: Logistics and Transportation Review, 2020, 135: 101866.\par
\bibitem{ref:sahin} Şahin M K, Yaman H. A branch and price algorithm for the heterogeneous fleet multi-depot multi-trip vehicle routing problem with time windows[J]. Transportation Science, 2022, 56(6): 1636--1657.\par
\bibitem{ref:12} Froger A, Jabali O, Mendoza J E, Laporte G. The electric vehicle routing problem with capacitated charging stations[J]. Transportation Science, 2022, 56(2): 460--482.
\end{thebibliography}
\end{document}
